\section{A Julia planning interface must preserve a single decision experiment}
\label{exec23:pomdpsjl-scope}

A bank considering Julia for sequential servicing needs to decide which part
of a planning system it is buying into. POMDPs.jl supplies a common vocabulary
for a model, solver, policy, belief updater and simulator. These interfaces are
valuable because the same declared problem can be presented to different
algorithms, but the interface itself neither identifies that problem from
observations nor proves the returned policy optimal. The root package at the
selected 8 June 2026 repository revision declares version 1.0.0. Sixteen
authenticated source and documentation files support the present inspection.
The root MIT license explicitly warns that other packages in the organization
can have different licenses. POMDPTools and individual solvers therefore need
their own source, rights and numerical reviews. Section~\ref{sec:exec24-policies-other}
develops that review for selected POMDPTools evaluation and simulation delegates;
its source findings do not constitute native qualification or coverage of all
solvers. The organizational name cannot stand in for a concrete adoption decision.
\citep{exec23_pomdpsjl_sources}

The model contracts distinguish a state space, an action space, a transition
distribution, an observation distribution, reward and discount. Let
$S_t$ be the account state before decision $A_t$, $S_{t+1}$ the resulting state,
$O_{t+1}$ the information released after that decision and $R_t$ its reward.
For a declared initial belief $b$ and policy $\pi$, a finite-horizon target is
\[
 V_H^\pi(b)=E_{b,\pi}\!\left[\sum_{t=0}^{H-1}\gamma^tR_t\right].
\]
This expression defines our supplied-model target. A causal interpretation
also needs a justified intervention law for $A_t$, sufficient state and an
honest observation clock. Merely writing a transition method does not supply
those premises. Moreover, the dynamic decision network used internally to
generate variables is a computational dependency graph; it is not a graph
discovered from bank data or an identified causal equivalence class.
\citep{exec23_pomdpsjl_sources,exec23_pomdpsjl_reference}

\subsection{Joint generation is an event, not three independent API calls}
\label{exec23:pomdpsjl-generation}

POMDPs.jl permits explicit distributions as well as a joint generator returning
a named tuple with next state, observation and reward. The documented caller
uses the \texttt{@gen} interface. Its inspected implementation invokes the
joint generator once for that invocation, takes supplied fields and fills
missing fields using transition, observation and reward methods in dependency
order. A POMDP observation can depend on both the old and new state, and the
reward fallback can depend on the observation as well. An adapter that calls
the generator separately for each output can change the joint experiment even
if every output retains the correct marginal distribution.
\citep{exec23_pomdpsjl_sources}

Consider a stipulated one-step distress event $U\sim\mathrm{Bernoulli}(1/4)$.
Set $S_{t+1}=U$, release alert $O_{t+1}=U$, and return reward $-40$ in distress
or $8$ otherwise. A coherent event has expected reward $-4$, always agrees
between state and alert, and has probability $1/4$ of an alert accompanying a
loss. If a wrapper independently regenerates the state, alert and reward,
the marginals survive, but state and alert agree with probability only $5/8$
and alert with loss occurs with probability $1/16$. The difference is material
to a policy: a subsequent rescue worth $35$ in distress and $-15$ otherwise,
taken only on an alert, has expected incremental value $35/4$ under the coherent
law and $-5/8$ under the split law. These are exact authored counterexamples
to a conversion strategy, not observed failures of the Julia package.
\citep{exec23_pomdpsjl_reference}

The proposed boundary therefore transports one event identifier and one
joint transition result, with model revision, state/action identifiers,
observation release time and reward units. A random seed alone does not define
equivalence across wrappers: the requested fields and invocation order also
determine how a joint generator consumes its random stream. Reproducibility
requires the concrete random-number implementation, serialized state or seed
protocol, call schedule and worker environment. The source exposes an explicit
random-number argument and a default when omitted; it does not provide the
bank's complete replay protocol.
\citep{exec23_pomdpsjl_sources}

\subsection{Observation, reward and termination need compatible signatures}
\label{exec23:pomdpsjl-contracts}

An observation method depending only on action and next state automatically
supports the fuller old-state signature through fallback. The reverse need
not hold. Suppose two accounts both reach state zero after the same action,
but an alert reports whether the previous state was one. Then the two histories
require alert probabilities zero and one, although action and next state are
identical. A solver requesting only those latter arguments cannot represent
both laws. The documentation explicitly acknowledges that some solvers require
the narrower signature. The repair is to select a compatible solver, augment
the modeled state with a justified sufficient history, or refuse the conversion;
silently dropping the old state changes the model.
\citep{exec23_pomdpsjl_sources,exec23_pomdpsjl_reference}

Reward conversion has a similar obligation. A solver asking for
$\bar r(s,a)$ needs the appropriate conditional expectation of a transition
reward, including the observation when relevant:
\[
 \bar r(s,a)=\sum_{s',o}T(s'\mid s,a)Z(o\mid s,a,s')r(s,a,s',o)
\]
in a finite explicit model. A joint generator can instead define the joint law
directly, in which case the expectation must follow that law. The inspected
documentation requires consistency between reward signatures. For the
two-outcome example, weighting by the declared probabilities gives $-4$;
averaging the two reward values uniformly gives $-16$. Returning a convenient
number from the narrower method is not enough.
\citep{exec23_pomdpsjl_sources,exec23_pomdpsjl_reference}

Terminal value is zero after termination in the model contract. That does not
erase the reward on the transition into the terminal state. A settlement
paying $12$ on entry has value $12$ with zero subsequent continuation; a wrapper
that masks the transition reward by destination terminal status reports zero.
This parallels the episode-clock issue in other sequential packages, without
asserting that the inspected Julia simulator implements that erroneous mask.
The root default marks no state terminal. A concrete bank model and simulator
must therefore bind terminal behavior, horizon, per-step discount and any
variable-duration convention explicitly. The inspected standard generator
fields do not by themselves establish a semi-Markov discount contract.
\citep{exec23_pomdpsjl_sources,exec23_pomdpsjl_reference}

\subsection{Belief meaning and a conditional Julia worker route}
\label{exec23:pomdpsjl-decision}

The library deliberately uses \emph{belief} broadly: it can be a posterior,
particles, a partial history or another policy information state. Its
initialization contract allows lossy conversion and has a pass-through default;
the history contract permits a partial suffix, and the default current
observation reads the last history item. Thus a successful initialization call
does not certify posterior normalization, sufficiency, nonempty history or an
independent snapshot. A wrapper must identify the representation and supported
operations, copy or serialize mutable state when isolation is required, and
refuse claims that exceed the retained information. State/action/observation
objects also require stable equality and identity: the model guide warns
against mutating objects used as dictionary keys or retained in histories.
\citep{exec23_pomdpsjl_sources}

The present disposition is a reference for interface design, with a conditional
future Julia worker. Full adoption is premature because an abstract
\texttt{solve} or \texttt{simulate} declaration is not an audited solver or
an executed experiment. A blanket port to R, Python, Java or C++ is also
unjustified: it would duplicate dispatch, distribution and stochastic-event
semantics before demonstrating a benefit. An independent finite-state oracle
remains useful precisely because it checks the intended experiment separately
from the chosen implementation.
\citep{exec23_pomdpsjl_sources,exec23_pomdpsjl_reference}

A prototype can proceed once it pins Julia and all dependency artifacts,
selects and reviews one concrete solver/updater/simulator combination, and
registers conformance cases for joint outputs, asymmetric transitions,
old-state observations, reward expectation, terminal entry, history loss,
state mutation and randomness. Measure offline preparation and online action
latency separately, along with compilation, serialization, memory, recovery
and numerical error at the required accuracy. The same target, information
clock, diagnostics and refusal must survive library, CLI and MCP use.
An LLM may help draft a model or map a business question into these contracts,
but a plausible generated method cannot authenticate a latent state or its
intervention law. Native integration, scientific qualification and institutional
adoption remain open.
\citep{exec23_pomdpsjl_sources}
